\subsection{Related work and the gap addressed}\label{sec:related}
\paragraph{Natural formation flight and ideal aerodynamics} The energetic benefit of flying in the upwash of a leader has been argued since early analyses of bird flight \cite{lissaman1970,hummel1983} and was later supported by physiological measurements on pelicans following a leader \cite{weimerskirch2001}. For aircraft, the first-order result is that a follower whose lifting line lies in a region of upwash needs less induced power at the same lift; the ideal limit for a pair of elliptically loaded wings and its dependence on lateral overlap, vertical offset and wake roll-up were worked out in the ideal-aerodynamics literature \cite{king2005}. Those results concern straight, steady flight, usually at high lift coefficient, where induced drag is a large share of total drag. A small survey UAV cruising at $\Va=\SI{18}{m/s}$ ($C_L\approx0.22$) sits at the opposite end of that range: induced drag is only \ncdiShare\,\% of total drag in our model, so even a perfect wake exploit can recover only a fraction of it.

\paragraph{Numerical models of close formation} Close-formation studies have used lifting-line and vortex-lattice models, panel codes and RANS CFD. Vortex-lattice models are cheap and capture the spanwise upwash gradient but, unless a free wake is used, assume a planar sheet \cite{katz2001,weissinger1947}; analytic models (a Lamb--Oseen pair \cite{lamb1932}) assume the roll-up is complete. RANS CFD of fixed-wing UAVs in formation has been reported to give large lift-to-drag increases at favourable offsets \cite{oendo2025} and to show that rolling moments vary strongly with offset. These studies are aerodynamic only: they contain no battery, no turns, no camera footprint and no controller, which is the gap that the mission-level coupling in this paper addresses.

\paragraph{Formation control and safety} Leader--follower formation control of fixed-wing aircraft is typically implemented with guidance-level predictive controllers around an autopilot \cite{beard2012}; collision avoidance is increasingly implemented as a safety filter built on control barrier functions \cite{ames2017} and their high-order extension for constraints of relative degree larger than one \cite{xiao2022}. Because a pointwise-in-time, fixed-gain radial HOCBF-QP enforces forward invariance only while a feasible input exists inside actuator limits, severe longitudinal closure can defeat a myopic radial filter unless input-constrained extensions---such as adaptive penalty dynamics (AdaCBFs) \cite{xiao2021ada}, finite-horizon lateral evasive backup trajectories (backup CBFs) \cite{chen2021backup}, or multi-axis fixed-wing CBFs \cite{molnar2025}---are employed; we evaluate the baseline radial filter explicitly under deceleration saturation and communication dropout (Sec.~\ref{sec:control}).

\paragraph{Cooperative mapping} Coverage planning for photogrammetry is usually posed as a boustrophedon problem in which energy is proportional to path length or time. A fixed-wing aircraft cannot turn between adjacent strips with a semicircle when the strip spacing is smaller than the turning diameter; the shortest curvature-bounded paths are Dubins paths \cite{dubins1957}, and for the spacing used here they are roughly twice as long as a semicircle (Sec.~\ref{sec:framework}). None of the works we are aware of couples this planning layer to a wake model; the present study does, but only with simulation-level fidelity.

\paragraph{Gap} We therefore ask a narrower and testable question than ``is formation flight beneficial'': \emph{under which transit distance, speed, manoeuvre cost, station-keeping accuracy and overlap assumptions does wake-aware formation reduce fleet energy per accepted hectare for a three-aircraft survey, and how much of any observed benefit is aerodynamic?}

\subsection{Research questions and hypotheses}
Table~\ref{tab:rq} restates the research questions of the preliminary version in a form that can be answered with the evidence actually collected here, and records the status of each hypothesis.
\begin{table}[H]\centering\small
\caption{Research questions and hypotheses with their status after this study.}\label{tab:rq}
\begin{tabular}{@{}p{0.5cm}p{5.0cm}p{3.1cm}p{5.0cm}@{}}\toprule
ID & Question / hypothesis & Evidence used & Status\\\midrule
RQ1 & How do offset and airspeed set the follower's induced-drag change, rolling moment, trim drag and electrical power? & strip theory, VLM, free wake, RANS CFD & Answered at model level; magnitude differs between models (spread reported)\\
RQ2 & Can NMPC + barrier-function filter keep separation and tracking under wind, delay and manoeuvres? & 540 guidance-level runs, step tests, ablation & Answered in simulation; myopic radial filter is conditional, fails under longitudinal deceleration saturation\\
RQ3 & Does a wake-aware adaptive formation reduce fleet energy per accepted hectare versus independent flight? & 200 paired mission trials $\times$ 5 models & Reduction found (\nabsASheet\,\%) but attributable to route/overlap planning, not to the wake\\
RQ4 & Does the formation/stereo geometry preserve photogrammetric quality (RMSE$_{H,V}$)? & geometry only & \textbf{Not answered} (no reconstruction)\\
RQ5 & Which factors govern the break-even regime? & Monte-Carlo sweeps, SRC, Sobol, tornado & Answered: distance, airspeed and wake effectiveness dominate\\\midrule
H1 & A follower at $d_y/b\in[0.75,0.90]$ draws $\ge4$\,\% less electrical power than isolated flight after trim drag & models, CFD & Supported by the broad-benefit models (\npeakFree--\npeakSheet\,\%), not by a sharp optimum at $0.80b$\\
H2 & Adaptive formation uses less energy than a rigid V & mission MC & Supported ($\nrelDSheet$\,\%)\\
H3 & Wake-aware adaptive formation lowers $\Earea$ by $\ge5$\,\% vs independent UAVs at equal quality & mission MC & Not supported ($\nrelASheet$\,\%, of which the wake contributes $\nrelBSheet$\,\% against)\\
H4 & Speed, distance, reconfiguration time and wake uncertainty govern break-even & SRC/Sobol & Supported for distance, speed, wake effectiveness; reconfiguration window is minor\\\bottomrule
\end{tabular}\end{table}

\subsection{Organisation of the paper}
Section~\ref{sec:model} presents the physical models, from the isolated aircraft to the OpenFOAM set-up; Sec.~\ref{sec:framework} the formation controller, supervisory logic and mission/coverage accounting; Sec.~\ref{sec:methods} the experimental design and verification; Sec.~5 the results (aerodynamics, CFD, control, reconfiguration, mission, break-even, coverage); Sec.~6 discussion, threats to validity and limitations. Appendices give the audit of the preliminary version, full parameter and scenario tables, additional statistics and the reproducibility protocol.
